\chapter*{Abstrakt}

Geschäftskritische, verteilte Legacy-Systeme werden häufig über ihren ur\-sprüng\-lich geplanten Lebenszyklus hinaus betrieben und weisen dabei strukturelle Design Smells auf, die Skalierbarkeit und Wartbarkeit beeinträchtigen. Direkte Eingriffe in den Quellcode solcher Systeme sind aufgrund fehlender Dokumentation und mangelnder Expertise häufig mit hohem Risiko verbunden. Eine vorangegangene Arbeit adressierte dieses Problem für ein reales Alarmempfangssystem mittels des \gls{lps}, das synchrone Kommunikationsketten durch eine transparente, \gls{mqtt}-basierte Proxy-Broker-Proxy-Architektur asynchron entkoppelt, ohne den bestehenden Quellcode zu verändern. Eine experimentelle Analyse im Rahmen dieser Masterarbeit zeigt jedoch, dass das \gls{lps} unter realistischen Fehlerbedingungen (Netzwerkfehler, Serverfehler) zentrale Delivery-Garantien verletzt: In den durchgeführten Experimenten traten bei $5\,000$ Alarmmeldungen $77{,}8\,\%$ Duplikate und $13{,}5\,\%$ Reihenfolgeverletzungen auf. Zusätzlich zeigt sich, dass die Skalierbarkeit des \gls{lps} primär durch das Hinzufügen weiterer Service-Instanzen erreicht wird, was den operativen Aufwand erhöht.

Aufbauend auf diesen Erkenntnissen entwickelt diese Masterarbeit das \gls{alps}, eine erweiterte Architektur, die auf \gls{amqp}-basierten Broker-Mechanismen wie Consistent-Hash-Routing, \texttt{x-single-active-consumer} und strikter Flusskontrolle (\texttt{prefetch\_count = 1}) beruht, um eine deterministische \textit{In-Order-Delivery} pro Alarmanlage sicherzustellen. Ergänzt wird die Architektur um eine zu\-sätz\-li\-che, Mnesia-gestützte Deduplizierungskomponente, die als vorgelagerte Sidecar-Komponente fachliche Duplikate erkennt und damit \textit{Exactly-Once-Delivery} Ende-zu-Ende gewährleistet. Die experimentelle Evaluation unter identischen Fehlerbedingungen zeigt, dass \gls{alps} beide Delivery-Garantien vollständig einhält ($0{,}0\,\%$ Duplikate, $0{,}0\,\%$ Reihenfolgeverletzungen) und zu\-sätz\-lich durch Queue-Partitionierung eine effizientere, operativ weniger aufwändige Skalierung ermöglicht als das ursprüngliche \gls{lps}. Die Ergebnisse belegen, dass \gls{alps} eine robuste und skalierbare Lösung für das Architektur-Reengineering geschäftskritischer, verteilter Legacy-Systeme darstellt, ohne Änderungen am bestehenden Quellcode zu erfordern.
